\clearpage\pagestyle{otMain}\setlength{\otBlocEspacement}{2.5mm}\fontsize{11}{15}\selectfont\corrigesChapitre{4}{Carrés et racines carrées}
\begin{blocLivre}{cor-4-1}\begin{corrige}{1}{Les carrés connus}\demarche{Relier chaque carré à un produit.}
$7\times7=\res{49}$, $11\times11=\res{121}$ et $12\times12=\res{144}$. Ces carrés sont des repères à mémoriser pour encadrer ensuite des racines.
\end{corrige}\end{blocLivre}
\begin{blocLivre}{cor-4-2}\begin{corrige}{2}{Lire trois racines}\demarche{Chercher le nombre positif dont le carré convient.}
$8^2=64$, $10^2=100$ et $0^2=0$. Les racines demandées sont donc $\res8$, $\res{10}$ et $\res0$, toutes positives ou nulles.
\end{corrige}\end{blocLivre}
\begin{blocLivre}{cor-4-3}\begin{corrige}{3}{Un signe à expliquer}\demarche{Utiliser la définition et non seulement le carré.}
Même si $(-9)^2=81$, une racine carrée est par définition positive ou nulle. On écrit donc $\res{\sqrt{81}=9}$.
\end{corrige}\end{blocLivre}
\begin{blocLivre}{cor-4-4}\begin{corrige}{4}{Encadrer entre deux entiers}\demarche{Placer le nombre entre deux carrés parfaits.}
$7^2=49<53<64=8^2$, donc $\res{7<\sqrt{53}<8}$. Aucun calcul approché n’est nécessaire.
\end{corrige}\end{blocLivre}
\begin{blocLivre}{cor-4-5}\begin{corrige}{5}{De l’aire au côté}\demarche{Inverser le calcul de l’aire du carré.}
Si le côté est $c$, son aire vaut $c^2$. La longueur cherchée est donc $\sqrt{121}=\res{11\ \mathrm{cm}}$. L’unité passe de cm$^2$ à cm.
\label{fig-cor-4-5}\figureGeo[corrige]{carre-121}
\end{corrige}\end{blocLivre}
\begin{blocLivre}{cor-4-6}\begin{corrige}{6}{Une valeur approchée}\demarche{Vérifier l’arrondi avec un encadrement.}
La calculatrice donne $4{,}358898\ldots$, donc $\sqrt{19}\approx\res{4{,}36}$. De plus $4{,}3^2=18{,}49<19<19{,}36=4{,}4^2$, d’où $\res{4{,}3<\sqrt{19}<4{,}4}$.
\end{corrige}\end{blocLivre}
\begin{blocLivre}{cor-4-7}\begin{corrige}{7}{Périmètre exact}\demarche{Arrondir seulement le résultat final.}
Le côté vaut $\sqrt{18}$ m, donc le périmètre exact est $\res{4\sqrt{18}\ \mathrm m}$. On calcule ensuite $4\sqrt{18}\approx\res{16{,}97\ \mathrm m}$. Arrondir le côté avant de le multiplier ferait perdre de la précision.
\label{fig-cor-4-7}\figureGeo[corrige]{carre-18}
\end{corrige}\end{blocLivre}
\begin{blocLivre}{cor-4-8}\begin{corrige}{8}{Deux aires}\demarche{Comparer les carrés puis leurs racines positives.}
$42<45$, donc $\sqrt{42}<\sqrt{45}$. Les deux côtés sont entre 6 et 7 cm car $36<42<45<49$. Leur différence est ainsi strictement comprise entre $\res0$ et $\res1$ cm.
\end{corrige}\end{blocLivre}
\begin{blocLivre}{cor-4-9}\begin{corrige}{9}{Détecter une fausse règle}\demarche{Calculer chaque membre indépendamment.}
Le membre de gauche vaut $\sqrt{100}=\res{10}$. Celui de droite vaut $6+8=\res{14}$. L’égalité est fausse : on ne peut pas répartir la racine sur une somme. Un seul contre-exemple suffit.
\end{corrige}\end{blocLivre}
\begin{blocLivre}{cor-4-10}\begin{corrige}{10}{Arrondi ou troncature}\demarche{Distinguer égalité exacte et approximation.}
$7{,}07^2=49{,}9849$, différent de 50 : le signe d’égalité est faux. $\sqrt{50}=7{,}071067\ldots$, donc $\res{\sqrt{50}\approx7{,}07}$ est bien l’arrondi au centième ; ici la troncature donne la même valeur.
\end{corrige}\end{blocLivre}
\begin{blocLivre}{cor-4-11}\begin{corrige}{11}{Un jardin carré}\demarche{Calculer le périmètre avant de retirer l’ouverture.}
Le côté vaut $\sqrt{72}$ m. La longueur à clôturer est donc $4\sqrt{72}-2\approx\res{31{,}9\ \mathrm m}$. On soustrait l’ouverture une seule fois au périmètre.
\label{fig-cor-4-11}\figureGeo[corrige]{jardin-72}
\end{corrige}\end{blocLivre}
\begin{blocLivre}{cor-4-12}\begin{corrige}{12}{Passer par une conversion}\demarche{Extraire une longueur puis convertir son unité.}
$\sqrt{6\,400}=\res{80}$ cm, soit $\res{0{,}8}$ m. On peut contrôler : $0{,}8^2=0{,}64$ m$^2$, ce qui correspond à 6 400 cm$^2$. Une conversion d’aire ne se fait pas avec le même facteur qu’une longueur.
\end{corrige}\end{blocLivre}
\begin{blocLivre}{cor-4-13}\begin{corrige}{13}{Une aire multipliée}\demarche{Revenir à l’aire avant de prendre la racine.}
L’aire initiale vaut 9 cm$^2$, la nouvelle 45 cm$^2$. Le côté vaut $\res{\sqrt{45}\ \mathrm{cm}}$, soit environ $\res{6{,}7\ \mathrm{cm}}$ au millimètre. Multiplier l’aire par 5 ne multiplie pas le côté par 5.
\label{fig-cor-4-13}\figureGeo[corrige]{carre-45}
\end{corrige}\end{blocLivre}
\begin{blocLivre}{cor-4-14}\begin{corrige}{14}{Un encadrement précis}\demarche{Ne pas confondre encadrement et arrondi.}
$3{,}16^2=9{,}9856$ et $3{,}17^2=10{,}0489$. Le nombre 10 est entre ces deux carrés, donc $\res{3{,}16<\sqrt{10}<3{,}17}$. Cet encadrement ne suffit pas seul à choisir l’arrondi au centième : il faudrait situer la racine par rapport à 3,165.
\end{corrige}\end{blocLivre}
\begin{blocLivre}{cor-4-15}\begin{corrige}{15}{Une plaque dans une boîte}\demarche{Comparer des carrés pour une décision exacte.}
Le côté de la plaque vaut $\sqrt{50}$ cm. Comme $7^2=49<50$, il dépasse 7 cm : la plaque ne rentre pas dans cette position. L’aire maximale d’une plaque carrée de côté 7 cm est $\res{49\ \mathrm{cm}^2}$. Une approximation grossière à 7 aurait masqué le dépassement.
\label{fig-cor-4-15}\figureGeo[corrige]{plaque-50}
\end{corrige}\end{blocLivre}
